% Direct-warp appendix shared by the modular and composed paper builds.
\providecommand{\WarpFigureRoot}{../support_geometry}
\appendices
\section{Direct Joint CONV Warps}
\label{app:directwarp}

The continuous CONV representation admits a direct extension from lattice
resize to rotation, skew, anisotropic affine transfer, and perspective warp.
The necessary change is not a decomposition of the coordinate map into
one-dimensional resampling passes.  It is a joint admission of the
two-dimensional current before that map is evaluated.  The construction
below consequently retains one inverse map, one continuous source atlas, and
one positive raster functional.  No intermediate raster occurs.

\subsection{One inverse map and its induced geometry}

Let \(\Omega\subset\R^2\) denote the normalized source rectangle and let
\(\widehat\Omega\subset\R^2\) denote the target rectangle.  A nonsingular
inverse warp
\begin{equation}
 F:\widehat\Omega\longrightarrow\Omega,\qquad q\longmapsto p=F(q),
 \label{eq:warpmap}
\end{equation}
is declared once.  Given the admitted source atlas \(U_Y\), point synthesis is
\begin{equation}
 (\mathcal W_FY)(q)=U_Y(F(q)).
 \label{eq:warppoint}
\end{equation}
Thus a shear factorization of an affine map is neither required nor used.
Indeed, repeated admitted interpolation would generally define a different
nonlinear operator, because every intermediate raster changes the current
state from which the next admission is formed.

When the source current geometry is expressed by the SPD field \(M(p)\), the
same inverse map induces the target quadratic form
\begin{equation}
 G_F(q)=J_F(q)^TM(F(q))J_F(q),
 \qquad J_F=DF,
 \label{eq:warppullback}
\end{equation}
and its determinant-one shape is
\begin{equation}
 \overline M_F(q)=\frac{G_F(q)}{\det(G_F(q))^{1/2}}.
 \label{eq:warpshape}
\end{equation}
Equation~\eqref{eq:warppullback} combines sampling scale, anisotropy, and the
warp before a square root or directional support is formed.  Consequently it
does not assume that Cartesian spacing commutes with the eigenframe of
\(M\).  The scalar \(\det(G_F)^{1/4}\) carries linear scale, while
\eqref{eq:warpshape} carries directional shape.

\subsection{Compiled bivariate proposal}

Index the source cells by \(C\), and write \((u,v)\in[0,1]^2\) for local
cell coordinates.  The joint source potential on cell \(C\) is represented
in the tensor Bernstein basis as
\begin{equation}
 U_C(u,v)=\sum_{i=0}^{5}\sum_{j=0}^{5}
 P_{C,ji}B_i^5(u)B_j^5(v).
 \label{eq:warppatch}
\end{equation}
The proposal controls are not fitted by an optimizer.  Let
\begin{equation}
 A_{ri}=B_i^5(r/5),\qquad 0\leq r,i\leq5,
 \label{eq:warpcollocation}
\end{equation}
and let \(Z_C\) contain the values of the complete two-order CONV proposal at
the \(6\times6\) fifth-step tensor lattice of \(C\).  Since the Bernstein
collocation matrix \(A\) is nonsingular, the unique bidegree-\((5,5)\)
proposal is
\begin{equation}
 \widetilde P_C=A^{-1}Z_CA^{-T}.
 \label{eq:warpproposal}
\end{equation}
All entries of every \(Z_C\) are obtained from one global fifth-step
evaluation.  Adjacent cells therefore have the same six samples on their
common edge; uniqueness in \eqref{eq:warpcollocation} then gives the same
complete Bernstein edge.  Equivalently, the patches are stored in one shared
global control lattice.  The resulting proposal is \(C^0\), and its current
is the gradient of one continuous piecewise-polynomial potential.

Let \(B\) denote the degree-elevated \(Q_1\) interpolant of the source
samples on that shared lattice.  It is cardinal, continuous, and affine
exact.  Admission will move \(B\) toward \(\widetilde P\), while the source
vertices remain fixed.

\subsection{Support-range admission}

For each source cell \(C\), let \(\mathcal S_C\) be the complete Cartesian
two-jet sample support, clipped at the outer boundary.  If shared control
entity \(e\) is incident to the cell set \(\mathcal I(e)\), define, for
channel \(c\),
\begin{equation}
 I_e^{(c)}=
 \bigcap_{C\in\mathcal I(e)}
 \left[
  \min_{p\in\mathcal S_C}Y_p^{(c)},
  \max_{p\in\mathcal S_C}Y_p^{(c)}
 \right].
 \label{eq:warprangeintersection}
\end{equation}
This intersection is nonempty: the corresponding \(Q_1\) control is a
convex combination of common incident vertices and belongs to every interval
in \eqref{eq:warprangeintersection}.  The proposal is first replaced by
\begin{equation}
 \widehat P_e^{(c)}=\Pi_{I_e^{(c)}}\widetilde P_e^{(c)}.
 \label{eq:warprangeclip}
\end{equation}
Unlike a post-reconstruction clamp, \eqref{eq:warprangeclip} constrains the
continuous potential.  Subsequent current admission remains on the segment
from \(B_e^{(c)}\) to \(\widehat P_e^{(c)}\), so every final control remains
in \(I_e^{(c)}\).

\begin{theorem}[Continuous support-range admission]
For every channel and source cell,
\begin{equation}
 \min_{p\in\mathcal S_C}Y_p^{(c)}
 \leq U_C^{(c)}(u,v)\leq
 \max_{p\in\mathcal S_C}Y_p^{(c)},
 \qquad (u,v)\in[0,1]^2.
 \label{eq:warprangetheorem}
\end{equation}
Consequently, point synthesis and every positive warped-pixel average have no
range excursion beyond the union of the source supports they traverse.
\end{theorem}

\begin{proof}
For cell \(C\), every one of its 36 shared controls belongs to the interval
on the right of \eqref{eq:warprangetheorem}.  The tensor Bernstein basis is
nonnegative and forms a partition of unity, hence \(U_C\) is a convex
combination of those controls.  Point evaluation gives the first result, and
integration against a nonnegative normalized measure gives the second.
\end{proof}

\subsection{Support-current cone}

The one-dimensional theorem controls current along a declared factor.  A
rotated or projective query instead requires a joint source current law.  For
cell \(C\), collect the four corner gradients of every \(Q_1\) cell contained
in \(\mathcal S_C\):
\begin{equation}
 \mathcal G_C=
 \left\{\nabla Q_D(z):D\subset\mathcal S_C,
              \ z\in\operatorname{vert}(D)\right\}.
 \label{eq:warpgenerators}
\end{equation}
Zero vectors are removed.  If the remaining directions generate a proper
closed planar cone, let
\begin{equation}
 K_C=\operatorname{cone}\mathcal G_C;
 \label{eq:warpcone}
\end{equation}
if no nonzero current is observed, or if the observed currents positively
span the plane, set \(K_C=\R^2\).  Thus a clean interface gives a narrow
cone, curvature gives a wedge, and crossing currents remove a directional
restriction rather than forcing a false owner or normal.

Every proper planar cone occurring here has a finite half-space description
\begin{equation}
 K_C=\left\{z\in\R^2:\ell_{Cr}^Tz\geq0,
                       \ 1\leq r\leq R_C\right\},
 \qquad R_C\leq3.
 \label{eq:warpconedual}
\end{equation}
A wedge uses its two inward boundary normals.  A ray additionally requires a
forward inequality, while a line uses two opposing transverse inequalities.
The whole-plane case has \(R_C=0\).

For each \(r\), the directional derivative
\begin{equation}
 L_{Cr}(u,v)=\ell_{Cr}^T\nabla U_C(u,v)
 \label{eq:warpdirectional}
\end{equation}
is degree elevated to the bidegree-\((5,5)\) Bernstein basis.  Denote its 36
linear control functionals by \(a_{Crk}^TP\).  The finite admission
conditions are
\begin{equation}
 a_{Crk}^TP\geq0,\qquad
 1\leq r\leq R_C,\quad 1\leq k\leq36.
 \label{eq:warpfiniteconstraints}
\end{equation}
Nonnegativity of these Bernstein coefficients is a sufficient certificate
for \(L_{Cr}\geq0\) on the complete cell.

The degree-elevated \(Q_1\) baseline is a constructive feasible point.  Its
gradient is a convex combination of the four corner gradients of its source
cell; those gradients occur in \(\mathcal G_C\).  Hence
\begin{equation}
 \nabla B_C(u,v)\in K_C
 \label{eq:warpbaselinefeasible}
\end{equation}
everywhere, and every margin in \eqref{eq:warpfiniteconstraints} is
nonnegative.

\subsection{Finite topological-entity admission}

Assign each nonvertex control to exactly one shared-edge or cell-interior
entity \(g\).  With the range-admitted proposal from
\eqref{eq:warprangeclip}, write
\begin{equation}
 \widehat P=B+\sum_gd_g,
 \label{eq:warpgroups}
\end{equation}
where the entity corrections have disjoint control ownership.  For a global
constraint \(a_k^TP\geq0\), put
\begin{equation}
 m_k=a_k^TB,\qquad
 H_k=\sum_g\max(0,-a_k^Td_g).
 \label{eq:warpmargin}
\end{equation}
If \(\widehat P\) already satisfies all constraints, it is accepted
unchanged.  Otherwise each harmful entity is assigned the simultaneous
capacity
\begin{equation}
 \alpha_g=min\!\left\{1,
  \min_{k:a_k^Td_g<0,\ H_k>0}\frac{m_k}{H_k}\right\},
 \label{eq:warpcapacity}
\end{equation}
with an empty inner minimum interpreted as one.  The first admitted state is
\begin{equation}
 P^{(0)}=B+\sum_g\alpha_gd_g.
 \label{eq:warpfirststate}
\end{equation}
Finally, let \(R=\widehat P-P^{(0)}\) and take the maximal common ray
\begin{equation}
 \tau=\min\!\left\{1,
   \min_{k:a_k^TR<0}
   \frac{a_k^TP^{(0)}}{-a_k^TR}\right\},
 \qquad P=P^{(0)}+\tau R.
 \label{eq:warpray}
\end{equation}
Again, an empty inner minimum is one.  Equations
\eqref{eq:warpcapacity}--\eqref{eq:warpray} contain neither an iteration nor
an active-set search.

\begin{theorem}[Finite joint current admission]
The control lattice \(P\) obtained from
\eqref{eq:warpcapacity}--\eqref{eq:warpray} satisfies every constraint in
\eqref{eq:warpfiniteconstraints}.  Consequently,
\begin{equation}
 \nabla U_C(u,v)\in K_C,
 \qquad (u,v)\in[0,1]^2,
 \label{eq:warpgradientcone}
\end{equation}
for every source cell.  The construction is finite, order independent,
cardinal, \(C^0\), and affine exact.
\end{theorem}

\begin{proof}
Fix constraint \(k\).  Contributions with \(a_k^Td_g\geq0\) can be
discarded in a lower bound.  For every harmful entity,
\(\alpha_g\leq m_k/H_k\), and therefore
\begin{equation}
 a_k^TP^{(0)}
 \geq m_k-\frac{m_k}{H_k}
       \sum_g\max(0,-a_k^Td_g)
 =0.
 \label{eq:warpcapacityproof}
\end{equation}
Equation~\eqref{eq:warpray} is the exact first intersection of the remaining
affine ray with the finite half-space boundary, hence \(a_k^TP\geq0\) as
well.  Applied to all derivative coefficients, Bernstein positivity gives
\(\ell_{Cr}^T\nabla U_C\geq0\), which is equivalent to
\eqref{eq:warpgradientcone} through \eqref{eq:warpconedual}.

Source vertices are not assigned to a movable entity.  Shared-edge controls
are single variables, so cardinality and \(C^0\) continuity are retained.
For affine data, the proposal equals the degree-elevated affine baseline and
already satisfies the cone and range constraints; it is consequently passed
unchanged.
\end{proof}

The theorem also explains why the construction is not an edge detector.  It
does not select one current from \(\mathcal G_C\).  Instead it admits the
complete high-order proposal whenever all of its derivative coefficients lie
in the cone generated by the currents actually present in the declared
support.

\subsection{Transport through affine and projective maps}

Let \(\gamma:[a,b]\to\widehat\Omega\) be a \(C^1\) target curve.  At every
regular point contained in one source cell, the chain rule gives
\begin{equation}
 \frac{d}{ds}U_Y(F(\gamma(s)))=
 \bigl[J_F(\gamma(s))\gamma'(s)\bigr]^T
 \nabla U_Y(F(\gamma(s))).
 \label{eq:warpchain}
\end{equation}
Let
\begin{equation}
 K_C^*=\{t\in\R^2:t^Tz\geq0\ \text{for all }z\in K_C\}
 \label{eq:warpdual}
\end{equation}
be the dual current direction cone.

\begin{corollary}[Warped-current transport]
Suppose \(F(\gamma(s))\) crosses only finitely many source-cell boundaries
and
\begin{equation}
 J_F(\gamma(s))\gamma'(s)\in K_{C(s)}^*
 \label{eq:warptangentcondition}
\end{equation}
for almost every \(s\), where \(C(s)\) contains \(F(\gamma(s))\).  Then
\(U_Y\circ F\circ\gamma\) is nondecreasing on \([a,b]\).
\end{corollary}

\begin{proof}
Equations \eqref{eq:warpgradientcone}, \eqref{eq:warpdual}, and
\eqref{eq:warpchain} make the derivative nonnegative in every open patch
segment.  The atlas is continuous and piecewise \(C^1\); therefore its
composition with the curve is absolutely continuous after splitting at the
finite knot set.  The fundamental theorem gives
\begin{equation}
 U_Y(F(\gamma(b)))-U_Y(F(\gamma(a)))
 =\int_a^b\frac{d}{ds}U_Y(F(\gamma(s)))\,ds\geq0.
\end{equation}
The knot points have measure zero and require no derivative convention.
\end{proof}

Affine maps transport the admitted directions by one constant Jacobian.
Homographies transport them by the spatially varying Jacobian of the same
declared map.  In either case, no target-space owner, orientation, or second
interpolation operation is introduced.

\subsection{Conservative warped-pixel analysis}

Point synthesis is appropriate for enlargement.  For reduction, let
\(\{P_j\}\) be the endpoint-aligned target Voronoi basins.  The matched
raster functional is
\begin{equation}
 (\mathcal B_FY)_j=
 \frac{1}{|P_j|}\int_{P_j}U_Y(F(q))\,dq.
 \label{eq:warppixel}
\end{equation}
Since the basins partition the target domain,
\begin{equation}
 \sum_j|P_j|(\mathcal B_FY)_j
 =\int_{\cup_jP_j}U_Y(F(q))\,dq.
 \label{eq:warpconservation}
\end{equation}
Thus \eqref{eq:warppixel} is conservative relative to the admitted continuous
profile.  It is also a positive average and inherits
\eqref{eq:warprangetheorem}.

For affine \(F\), a target basin maps to a source parallelogram. Splitting
at the integer source-knot lines produces convex polygons, each contained in
one patch. Positive fan triangulation introduces no overlap. On every
triangle, \(U_C\) has total degree at most ten, so the positive 25-node
triangle rule of degree ten integrates it exactly in real arithmetic.
Axis-aligned rectangular footprints use separable integrated Bernstein
weights. These are the affine cases of the moment-bank construction in
Appendix~\ref{app:convstarwarp}.

For a homography, source-coordinate integration yields
\begin{equation}
 (\mathcal B_FY)_j=\frac1{|P_j|}
 \int_{F(P_j)}U_Y(p)\left|\det DF^{-1}(p)\right|\,dp.
 \label{eq:warpprojectiveintegral}
\end{equation}
With homogeneous coordinate normalizations absorbed into the nonsingular
matrix \(H\), the inverse-map density is
\begin{equation}
 \left|\det DF^{-1}(p)\right|
 =\frac{|\det H^{-1}|}{|\ell(p)|^3},
 \label{eq:warpdensity}
\end{equation}
where \(\ell\) is the affine denominator of \(H^{-1}\). On a clipped
triangle, its vertex extrema \(\ell_-,\ell_+\) determine
\begin{equation}
 c_T=\frac{\ell_-+\ell_+}{2},\qquad
 r_T=\frac{\ell_+-\ell_-}{2|c_T|}.
\end{equation}
The canonical construction selects a positive triangle rule from the fixed
geometric table in Section~\ref{sec:stargeometric}. Its 25--55 nodes evaluate
\eqref{eq:warpdensity} directly. A triangle whose radius exceeds the final
table entry is partitioned into four midpoint children. This construction
uses the same geometric partition for every source field on the given
lattice and map.

The resulting 36 moment weights are contracted with the admitted controls.
A constant reference \(P_0\) is integrated through exact target coverage,
and the numerical weights act on \(P_{C,ij}-P_0\). With
\(m_T=\int_T|\det DF^{-1}|\) and
\(D_T=\max_{i,j}|P_{C,ij}-P_0|\), the fixed table gives
\begin{equation}
 E_T\leq \varepsilon_gD_Tm_T,\qquad \varepsilon_g=10^{-9},
 \label{eq:warpprojectivecertificate}
\end{equation}
in exact arithmetic. The geometric moment theorem proves this inequality
for all control values. Summing the triangle contributions and dividing by
the target-basin area gives the pixel-functional bound. Constant fields
retain exact target coverage; all channels reuse the same moment bank.

\subsection{Multichannel form}

For \(Y:\Omega\cap\mathbb Z^2\to\R^C\), the control topology, support sets,
and inverse map are shared.  Range and current admission are applied to each
component in the declared channel basis, yielding the product set
\begin{equation}
 \prod_{c=1}^C
 \left\{U^{(c)}:\nabla U_C^{(c)}\in K_C^{(c)},\ 
 U_C^{(c)}\in I_C^{(c)}\right\}.
 \label{eq:warpmultichannel}
\end{equation}
The preceding proofs therefore hold componentwise.  As with the main CONV
operator, this statement is invariant under independent affine changes of
channel units; it does not assert invariance under an arbitrary basis change
that mixes channels.

\subsection{Exact-truth evaluation}

The evidence uses only procedural sources whose point values and warped-pixel
means are independent of every tested interpolator.  The sharp corpus
contains a half-plane with normal angle \(37^\circ\), a width-0.065 bar at
\(27^\circ\), two width-0.055 crossing bars at \(31^\circ\) and
\(121^\circ\), and a \(0.62\times0.34\) box at \(-19^\circ\), all centered
at \((1/2,1/2)\).  The four target-to-source maps are
\begin{align}
 F_{\rm rot}&: (\theta,s_x,s_y)=(21^\circ,0.72,0.72),\nonumber\\
 F_{\rm skew}&: (s_x,s_y,h_x)=(0.70,0.75,0.42),\nonumber\\
 F_{\rm aniso}&: (\theta,s_x,s_y)=(-14^\circ,0.50,0.84),\nonumber\\
 F_{\rm proj}&: (g_x,g_y,\theta,s_x,s_y)
 =(0.32,-0.18,9^\circ,0.67,0.72),
 \label{eq:warpcorpusmaps}
\end{align}
where affine maps use \(F(q)=\tfrac12+A(q-\tfrac12)\), with
\(A=R_\theta\left(\begin{smallmatrix}1&h_x\\h_y&1\end{smallmatrix}\right)
\operatorname{diag}(s_x,s_y)\), and the projective map uses
\begin{equation}
 F(q)=\frac12+
 \frac{A(q-\frac12)}{1+(g_x,g_y)^T(q-\frac12)}.
 \label{eq:warpprojectivedeclared}
\end{equation}
Exact polygon membership supplies point truth.  Since a source half-plane
remains a target half-plane under a pole-free homography, polygon clipping in
target coordinates supplies exact basin coverage.

The smooth census uses the real Fourier fields
\begin{equation}
 f(x)=b+\sum_{k=1}^Ka_k
 \cos\!\left(2\pi\xi_k^Tx+\varphi_k\right).
 \label{eq:warpfourierfield}
\end{equation}
For an affine map, the exact mean over a target rectangle with center \(m\)
and side lengths \(h_x,h_y\) is obtained by multiplying each transformed
mode by
\begin{equation}
 \operatorname{sinc}((A^T\xi_k)_xh_x)
 \operatorname{sinc}((A^T\xi_k)_yh_y).
 \label{eq:warpfouriermean}
\end{equation}
The scalar probes use the two mode sets
\begin{equation}
\begin{aligned}
 \mathcal F_1=\{&((5,2),.22,.13),\ ((9,3),.09,1.17),\\
                &((2,-4),.07,-.41)\},\\
 \mathcal F_2=\{&((6,1),.17,.2),\ ((-1,6),.17,-.6),\\
                &((4,4),.09,1.1)\},
\end{aligned}
\label{eq:warpfourierprobes}
\end{equation}
with \(b=0.5\).  A three-channel probe uses wavevectors
\((4,1),(1,5),(5,-3)\), phases \(0,0.7,-0.4\), and amplitude rows
\begin{equation}
 \begin{bmatrix}.22&.05&.08\\.04&.20&.07\\.06&.04&.18\end{bmatrix},
 \qquad b=(.5,.5,.5).
 \label{eq:warpmultichannelprobe}
\end{equation}

Table~\ref{tab:warpsharp} reports the 16 sharp cases on a
\(17\times17\) source and \(13\times13\) target.  CONV has the lowest MSE in
seven point cases and seven warped-pixel cases.  More importantly for the
proved admission, its maximum range excursion remains at floating roundoff,
whereas the factor-only ablation and Lanczos admit substantial excursion.

\begin{table}[H]
\caption{Sharp affine/projective exact-truth census}
\label{tab:warpsharp}
\centering
\scriptsize
\setlength{\tabcolsep}{2.0pt}
\begin{tabular}{lrrrr}
\hline
Method & Point MSE & Pixel MSE & Point exc. & Pixel exc.\\
\hline
CONV & .021002 & .007377 & \(2.22\!\times\!10^{-16}\) & \(3.18\!\times\!10^{-14}\)\\
Factor ablation & .021137 & .007705 & .2287 & .1433\\
Lanczos-3 & .022928 & .008741 & .2817 & .1989\\
Bilinear & .021513 & .007576 & 0 & 0\\
\hline
\end{tabular}
\end{table}

\begin{figure*}[!t]
 \centering
 \includegraphics[width=0.98\textwidth]{\WarpFigureRoot/conv_warp_projective_plate_canonical.png}
 \vspace{3pt}
 \includegraphics[width=0.74\textwidth]{\WarpFigureRoot/conv_warp_canonical_metrics.png}
 \caption{Direct joint CONV warp evidence.  Top: one homography applied to
 sources with exact projective basin truth; each displayed certificate is the
 computed enclosure of the projective polynomial integral.  Bottom:
 topology and exact-truth fidelity over the declared corpus.  Values plotted
 at the numerical floor are below \(2\times10^{-16}\).  The factor-only
 result is an ablation of joint current admission, not a second CONV
 definition.}
 \label{fig:warpevidence}
\end{figure*}

The intrinsic topology census evaluates normal profiles rather than using MSE
as a surrogate for ringing.  For a profile \(z_0,\ldots,z_n\) whose exact
orientation is increasing, the reported wrong-way current is
\begin{equation}
 V_-(z)=\sum_{i=0}^{n-1}\max(0,z_i-z_{i+1}).
 \label{eq:warpwrongway}
\end{equation}
Straight interfaces use six source angles, three maps, 513 samples per
profile, and source sides 17 and 33.  CONV has no normal sign changes,
\(V_-\leq1.34\times10^{-14}\), and excursion below
\(4.45\times10^{-16}\).  Lanczos-3 produces four to nine sign changes,
mean \(V_-=0.276\)--\(0.277\), and mean excursion \(0.166\).

For curved interfaces, circles and ellipses are tested at 12 phases under
three maps, giving 72 profiles per source size.  Table~\ref{tab:warpcurved}
shows that the admitted cone eliminates wrong-way normal current at sides 13,
17, and 25.  At side 9, the range theorem remains exact, while the six-node
support spans competing conic normals and the cone correctly widens; the
remaining wrong-way current therefore measures support ambiguity rather than
failure of the finite inequalities.

\begin{table}[H]
\caption{Curved-interface intrinsic normal census for CONV}
\label{tab:warpcurved}
\centering
\footnotesize
\setlength{\tabcolsep}{4pt}
\begin{tabular}{rrrr}
\hline
Source side & Mean \(V_-\) & Max. \(V_-\) & Max. excursion\\
\hline
9  & \(1.94\times10^{-3}\) & \(2.80\times10^{-2}\) & \(2.22\times10^{-16}\)\\
13 & \(1.67\times10^{-15}\) & \(4.77\times10^{-15}\) & \(4.44\times10^{-16}\)\\
17 & \(2.23\times10^{-15}\) & \(6.44\times10^{-15}\) & \(4.44\times10^{-16}\)\\
25 & \(3.31\times10^{-15}\) & \(5.44\times10^{-15}\) & \(4.44\times10^{-16}\)\\
\hline
\end{tabular}
\end{table}

On the nine band-admitted Fourier cases, mean point MSE is
\(1.48764\times10^{-4}\) for CONV, \(1.06820\times10^{-4}\) for the
factor-only ablation, \(4.57424\times10^{-5}\) for Lanczos-3, and
\(1.22893\times10^{-3}\) for bilinear.  The corresponding warped-pixel MSEs
are \(6.27668\times10^{-5}\), \(4.09234\times10^{-5}\),
\(1.88691\times10^{-5}\), and \(7.63065\times10^{-4}\).  Thus the complete
support maximum principle has a measurable smooth-field fidelity cost; the
experiment does not support universal MSE dominance.

Finally, affine conservation equates the weighted target sum and an
independent exact integral of the admitted atlas to \(3\times10^{-12}\).  On
a nonconstant projective plane, independent 40-point tensor integration
differs from the certified result by \(1.18\times10^{-11}\), within its
\(5.84\times10^{-10}\) enclosure.  The largest certificate returned by the
sharp projective corpus is \(1.52\times10^{-11}\).  These computations verify
the exact polynomial and error-enclosure implementations; the cardinality,
integrability, range, current-cone, and transport statements follow from the
preceding formulas rather than from the measurements.

The direct warp therefore extends CONV without changing its governing idea:
the high-order proposal is admitted in the representation that carries the
current, after which point evaluation and positive basin analysis are two
functionals of the same continuous object.  The principal remaining boundary
is source resolution itself.  When one compact two-jet support contains
genuinely competing interface normals, the observed cone contains both and
cannot certify a narrower direction without adding information absent from
the samples.
